\subsection{A conditional heat-bath criterion for the Wilson fibre}
\label{ssec:fibre-heatbath}

The orthogonal projection onto the constants in the internal-fibre
framework of James Montgomery Lockwood~\cite{Loc26Fibres} suggests an
analogous projection on a constrained lattice fibre. Here the measure is
the Wilson posterior, and each projection averages a single free link
conditional on all other free links. These projections act on a
finite-volume probability space. Their generator and its spectral gap
are different from the vertical Laplacian and its gap in the internal
scalar model.

\begin{lemma}[From fibre relaxation to covariance]\label{lem:fibre-heatbath}
Let $\nu$ have a positive density on a finite product and suppose its
single-coordinate conditional laws depend only on the neighbours in an
undirected graph of maximum degree at most $\Delta$, where $\Delta\ge1$.
Write $P_i$ for conditional expectation given all coordinates except $i$,
$Q_i=I-P_i$, and $L=\sum_i Q_i$. If
$\langle h,Lh\rangle_\nu\ge\lambda\operatorname{Var}_\nu(h)$ for every real
$h\in L^2(\nu)$, with $\lambda>0$, then bounded real $f,g$ supported on
coordinate sets $A,B$ at graph distance $r\ge1$ satisfy
\[
 |\operatorname{Cov}_\nu(f,g)|\le
 \|f\|_\infty\|g\|_\infty
 \left[\frac{4|A|}{e\Delta}e^{-r/2}
        +e^{-\lambda r/(2e\Delta)}\right].
\]
The right side is taken to be zero if there is no path from $A$ to $B$.
\end{lemma}

\begin{proof}
Let $\delta_j h$ be the supremum of $|h(u)-h(u')|$ over configurations
differing only at $j$. For $j\ne i$, subtract the two conditional
integrals defining $P_i h$, first at fixed conditional law and then at
fixed integrand. This gives
$\delta_j P_i h\le\delta_j h+\mathbf{1}_{i\sim j}\delta_i h$.
Also $\delta_i P_i h=0$. The Markov operator $I-sL$ is a convex
combination of the identity and the $P_i$ for $0\le s\le1/|I|$, where
$|I|$ is the number of coordinates. Thus
$\delta((I-sL)h)\le (I-sI+s\mathcal A)\delta(h)$ componentwise, with
$\mathcal A$ the graph adjacency matrix. Apply this to Euler products
and use uniform convergence of the bounded-operator exponential to get
$\delta(e^{-tL}h)\le e^{t(\mathcal A-I)}\delta(h)$.
For $j\in A$, the terms of order $k<r$ vanish when applied to $\delta(g)$,
and the total weight of walks of length $k$ from $j$ is at most
$\Delta^k$. Hence
\[
 \delta_j(e^{-tL}g)\le 2\|g\|_\infty e^{-t}
   \sum_{k\ge r}\frac{(t\Delta)^k}{k!}
 \le 2\|g\|_\infty\exp(-r+e\Delta t).
\]
Each $Q_i$ is an orthogonal projection, $Q_i f=0$ for $i\notin A$, and
$\|Q_i h\|_\infty\le\delta_i h$. Therefore
\[
 \big|\langle f,g-e^{-TL}g\rangle_\nu\big|
 \le \sum_{i\in A}\int_0^T
       \big|\langle Q_i f,Q_i e^{-tL}g\rangle_\nu\big|\,dt
 \le \frac{4|A|}{e\Delta}\|f\|_\infty\|g\|_\infty e^{-r+e\Delta T}.
\]
The spectral assumption bounds
$|\operatorname{Cov}_\nu(f,e^{-TL}g)|$ by
$e^{-\lambda T}\sqrt{\operatorname{Var}_\nu(f)\operatorname{Var}_\nu(g)}$.
Set $T=r/(2e\Delta)$. If $r=\infty$, the walk sum is zero for all
$t$ and one lets $T\to\infty$.
\end{proof}

\begin{corollary}[A sufficient relaxation rate for \textup{(FC)}]
\label{cor:fibre-heatbath}
For the straight-path Wilson fibres of Proposition~\ref{prop:bpfibreduce},
let $b\ge16$ be the blocking factor, and choose the marks nearest the
midpoints of the paths. Suppose the single-free-link heat-bath operators
have $L^2$ gaps at least $\lambda_b>0$, uniformly over all coarse data and
all torus sizes under consideration. Put $a_b=b\min(1,\lambda_b)$.
For a fixed matrix gauge group there are constants $C,c>0$, independent
of $b$, the coarse data and the torus, such that if $a_b\ge1$ then
\[
 \sup_j\sum_{e\ne j}J_{ej}(H^c)\le C\beta_b^2 e^{-c a_b}.
\]
The same sum over coarse separations at least $R\ge1$ is at most
$C\beta_b^2 e^{-c a_b R}$. In particular, a rate with
$a_b/\log(2+|\beta_b|)\to\infty$ gives the row and tail bounds in
\textup{(FC)} for all sufficiently large $b$.
\end{corollary}

\begin{proof}
In the marked gauge the remaining coordinates are the free fine links.
Join two such links if they belong to a common fine plaquette. The Wilson
action gives the required locality of each conditional law. Each fine link
belongs to at most six plaquettes and the graph has maximum degree at
most $18$. The force $D_X^a\mathcal S$ at a marked link $a$, for a unit
Lie-algebra vector $X$, is a function of the at most $18$ free links
sharing one of those plaquettes with $a$. Its supremum norm is at most
$C_G|\beta_b|$, since each term is a derivative of the trace of a
product of unitary matrices and each matrix has bounded norm.

Adjacent links in the plaquette graph have midpoint sup distance at most
one on the fine torus. If the starting vertices of the marked paths
have coarse torus sup distance $d\ge1$, their fine midpoint distance
is at least $bd-b/2$. If the starts coincide, distinct marks have
distinct orientations and midpoint distance at least $b/2-1/2$.
Passing from a mark to its support loses at most one in this distance
at each end. Thus the two force supports have graph distance at least
$(b/4)\max(1,d)$ for $b\ge16$. If a force has empty support, its covariance
is zero. Proposition~\ref{prop:bpfibreduce} identifies the mixed
derivative for distinct marks with minus the covariance of the two
forces. Apply Lemma~\ref{lem:fibre-heatbath} and take the supremum over
coarse data and unit tangent directions. This gives the entrywise bound
$C\beta_b^2\exp[-c a_b\max(1,d)]$.
The number of coarse links at sup distance $d$ is at most $C(1+d)^3$,
uniformly on tori, by counting lattice shells. Summing the estimate and
using $a_b\ge1$ gives both inequalities, after reducing $c$ to absorb the
polynomial shell factor.
\end{proof}

This criterion reduces a Wilson force bound to a rate of relaxation for
an explicit constrained heat bath. Finite-volume positivity alone gives
only a gap bound by the density ratio $\exp(-\operatorname{osc}\mathcal S)$.
Indeed, the heat-bath quadratic form is a sum of squared distances to
the single-link conditioning subspaces. Comparing each squared distance
and the variance with the corresponding product-Haar quantity proves
the density-ratio bound, since the product-Haar heat bath has gap one.
That comparison gives no volume-uniform rate along the tuning. A proof
of the hypothesis of Corollary~\ref{cor:fibre-heatbath} would still
require uniform control of the constrained Wilson measures, including
rough coarse data.
